\worksheettitle
  {Единицы длины}
  {\modulemark{M03} Версия учителя}

\begin{teacherbox}[Паспорт модуля]
  \textbf{Результат:} ученик переводит длины между основными единицами,
  сохраняет величину и объясняет преобразование через соотношение единиц.

  \textbf{Предварительно:} понимает разрядное устройство числа, знает
  умножение и деление на 10, 100 и 1000, различает отрезок и его длину.
\end{teacherbox}

\begin{teacherbox}[Методическое ограничение]
  Не сводите перевод к правилу «дописать нули». Перед вычислением ученик
  называет соотношение единиц и прогнозирует, увеличится или уменьшится число.
\end{teacherbox}

\section{Вход и примеры}

Во входе: $10$ мм, $100$ см, $1000$ м; сантиметр меньше метра.

\unitsladderfigure

В первом примере пропуск: $10$ мм. Проверка второго примера:
$9\cdot100+5=905$ см.

\section{Ответы к практике}

\begin{practicebox}[1. В одну единицу]
  a) $64$ мм; б) $38$ см; в) $508$ см; г) $6045$ м.
\end{practicebox}

\begin{practicebox}[2. Разложение]
  a) $24\text{ см }8\text{ мм}$; б) $7\text{ м }30\text{ см}$;
  в) $7\text{ км }300\text{ м}$.
\end{practicebox}

\begin{practicebox}[3. Самостоятельно]
  a) $47$ см; б) $270$ см; в) $4\text{ км }65\text{ м}$;
  г) $1\text{ м }20\text{ см }8\text{ мм}$.
\end{practicebox}

\begin{teacherbox}[Диагностика составной записи]
  В $1208$ мм сначала выделяется $1000$ мм, то есть $1$ м. Остаток
  $208$ мм — это $20$ см и $8$ мм. Просите выполнить обратную проверку.
\end{teacherbox}

\section{Ошибка и контроль}

\begin{warningbox}[Исправление]
  Причина: $6$ км ошибочно прочитаны как $600$ м или цифры просто
  соединены. Так как $1\text{ км}=1000\text{ м}$,
  \[
    6\text{ км }45\text{ м}=6000\text{ м}+45\text{ м}=6045\text{ м}.
  \]
\end{warningbox}

\begin{checkpointbox}[Ответы]
  \begin{enumerate}[label=\textbf{\arabic*.}]
    \item $4080$ м;
    \item $6\text{ м }5\text{ см}$;
    \item величины равны: $3\text{ м }8\text{ см}=308\text{ см}$;
    \item сантиметр в $100$ раз меньше метра, поэтому для той же длины
      требуется в $100$ раз больше сантиметров.
  \end{enumerate}
\end{checkpointbox}

\begin{teacherbox}[Решение о маршруте]
  M03-R назначается, если ученик не использует соотношение единиц, неверно
  прогнозирует направление изменения числа или соединяет цифры составной
  записи. M03-H достаточно при единичных вычислительных ошибках. Устойчивый
  результат открывает M04 и M07.
\end{teacherbox}
